%% ===================================================================
%% PART V
%% ===================================================================
\part{Relation to established mathematics, provenance, and open questions}

%% -------------------------------------------------------------------
\section{Relations to established mathematics}\label{sec:literature}

This section compares the objects and results of the paper with standard notions and theorems in neighbouring areas. It is a comparison of definitions and statements, not a survey of the literature, and \emph{no claim of novelty or priority is made}.

\paragraph{What was read, and how.} The following sources were read in full text (as retrieved on the date of the manuscript): Heden's survey of vector-space partitions \cite{Heden2012}; Kurz's paper on Heden's tail bound \cite{KurzTail} (for the exact statement of Heden's theorem); Kurz's classification of vector-space partitions of $\mathrm{GF}(2)^8$ \cite{KurzGF28}; Cooper--Okur on multigraphs with a unique partition into cycles \cite{CooperOkur}; Margolis--Saliola--Steinberg (for the statement and axioms of the Bohne--Dress theorem on zonotopal tilings) \cite{MSS}; and D.~S.~Kim's paper \cite{Kim2009} (for the statement of the Pless power-moment identities, which it attributes to MacWilliams--Sloane). The remaining references (MacWilliams--Sloane, Huffman--Pless, Oxley, Graver, Sturmfels, Schrijver, Ziegler, Matou\v sek, Hirschfeld, Beutelspacher, Terras, O'Donnell, Hemmecke) were assembled from catalogue records and author or publisher pages and were \emph{not} read; statements attributed to them are general background to be checked by the reader and no theorem number from them is quoted. Where a theorem of a read source is used below, it is stated as it appears there, and any consequence drawn is marked as a consequence derived here and not as a statement of the source.

\subsection{Binary codes}

The code $C=\ker M$ is an ordinary binary linear code \cite{MacWilliamsSloane1977,HuffmanPless2003}. A \emph{circuit} in our sense is a word of $C$ whose support is minimal under inclusion among supports of nonzero words; in coding theory such words are usually called \emph{minimal codewords}. Lemma~\ref{lem:circuit-basic}(b) is the familiar fact that every codeword is a disjoint union of minimal codewords (a fact that is special to the binary case). The reduced code, in which no two coordinates are series-equivalent and none is dead, is the code with no repeated or zero columns in its generator matrix; equivalently the dual code has minimum distance at least~$3$, i.e.\ $C$ is an orthogonal array of strength~$2$, and Lemma~\ref{lem:moment} is the corresponding second moment of the weight distribution. The family $\mathrm{RM}(1,k-1)$ of Proposition~\ref{prop:RM} is the first-order Reed--Muller code, a classical code with weights $0,p/2,p$. The even-weight code $E_p$ (parity-check code) is the simplest example of a code whose minimal codewords are pairs.

We did not find, in the sources read, the study of the \emph{set of partitions of the all-ones word into minimal codewords} as an object in its own right (the object of Part~III), nor the specific link between such partitions and vector-space partitions of $C/\langle\one\rangle$ (Theorem~\ref{thm:vsp}); the subsection on vector-space partitions below records precisely what \emph{is} classical in that part. This is an observation about the sources read, not a statement about the literature as a whole.

\subsection{Matroids and graphs}

The circuits of $C$ are the circuits of the binary matroid represented by the parity matrix $M$ (its column matroid), and equally the cocircuits of the vector matroid represented by a generator matrix; $C$ is the cocycle space of the latter \cite{Oxley2011}. The decomposition of a binary cocycle (or cycle) into disjoint circuits is standard for binary matroids; non-uniqueness of that decomposition is the phenomenon behind Theorem~\ref{thm:E1}.

\begin{example}[Cut spaces of graphs]\label{ex:cut}
If $C$ is the cut space of a connected graph $G$ (the set of edge cuts $\delta(S)$, as words in $\F_2^{E(G)}$), the circuits of $C$ are the \emph{bonds} (minimal nonempty cuts), and Theorem~\ref{thm:E1} reads: the box is non-embedded if and only if some cut has two distinct decompositions into edge-disjoint bonds. The even-weight code $E_p$ is the cut space of the cycle graph $C_p$ (bonds are pairs of edges), and the cut $\delta(\{1,3\})=\{12,23,34,41\}$ of $C_4$ has the three decompositions $\{12,34\}$, $\{23,41\}$ and so on; in this language $E_4$ is the smallest obstruction. The $K_{2,t}$ cut spaces mentioned in the release series (\gid{G10-K}, \CHECKED{} for $t\le10$ in the source) belong to this family; we do not treat them here.
\end{example}

\begin{remark}[Cycle codes and the theorem of Cooper and Okur]\label{rem:cooperokur}
For the \emph{cycle space} of a (multi)graph $G$, the circuits are the cycles, the codewords are the Eulerian subgraphs, and Theorem~\ref{thm:E1} says: the box is non-embedded iff some Eulerian subgraph has two different partitions of its edge set into cycles. Cooper and Okur \cite{CooperOkur} (Corollary~3.1 for multigraphs, Theorem~2.8 for digraphs, in the numbering of the arXiv version read) prove that a connected Eulerian multigraph has a \emph{unique} partition into cycles if and only if it is a bridgeless cactus multigraph, i.e.\ no two of its cycles share an edge. Combining the two statements (a translation made here; their paper does not use cycle spaces, matroids or codes): \emph{the box of the cycle code of a multigraph $G$ is embedded if and only if every connected Eulerian subgraph of $G$ is a cactus.} The smallest obstruction is two vertices joined by four parallel edges, whose cycle code is $E_4$ and whose three partitions into digons are the three partitions of $E_4$ in Section~\ref{sec:E4}. This graphic specialisation of Theorem~\ref{thm:E1} is thus a known theorem of graph theory, and our criterion extends it to arbitrary binary matroids. As a consistency check of the translation (not a proof), the equivalence between the criterion of Theorem~\ref{thm:E1} and ``some connected Eulerian subgraph contains two cycles sharing an edge'' was tested exhaustively on $2000$ random multigraphs with at most $5$ vertices and $9$ edges, with no discrepancy (\texttt{scripts/lit\_checks.py}). For cut spaces of planar graphs the duality (bonds of $G$ are cycles of the dual) transfers the same statement to bond partitions; we have not examined which of the minimal classes of Part~IV are graphic or cographic.
\end{remark}

\subsection{Finite affine geometry}

A spanning subset $X\subseteq\F_2^k$ determines the code $C(X)$, and the circuits of $C(X)$ are complements of hyperplanes of $X$.

\begin{lemma}[Circuits as hyperplane complements]\label{lem:hyp}
For $f\in\F_2^k\setminus\{0\}$ the word $w_f$ is a circuit of $C(X)$ if and only if $X\cap f^\perp$ spans the hyperplane $f^\perp$ (that is, has rank $k-1$).
\end{lemma}

\begin{proof}
For $g\ne0$, $\supp w_g\subseteq\supp w_f$ iff $X\cap f^\perp\subseteq X\cap g^\perp$. If $X\cap f^\perp$ has rank $k-1$ this forces $g^\perp\supseteq f^\perp$, so $g=f$. If the rank is smaller, choose $g\notin\{0,f\}$ vanishing on $\spn(X\cap f^\perp)$: then $w_g\ne w_f$ is a nonzero word with support inside that of $w_f$, hence strictly smaller (the supports cannot coincide because $f\mapsto w_f$ is injective).
\end{proof}

This is the test used in the source checker and in our engine (Appendix~\ref{app:repro}). Specialising to two-block partitions gives a geometric meaning for $n_2$.

\begin{proposition}[Two-block partitions as parallel hyperplane pairs]\label{prop:geo-n2}
Let $X\subseteq H\cong\AG(k-1,2)$ be a spanning subset and $\varphi$ a non-constant affine function on $H$. Then $\{X\cap\{\varphi=1\},\,X\cap\{\varphi=0\}\}$ is a partition of $\one$ into two circuits of $C(X)$ if and only if $X\cap\{\varphi=0\}$ and $X\cap\{\varphi=1\}$ each \emph{affinely span} the affine hyperplane $\{\varphi=0\}$, respectively $\{\varphi=1\}$, of $H$. Consequently
\[
\begin{aligned}n_2&=\#\bigl\{\text{directions of parallel hyperplane pairs of }\AG(k-1,2)\\&\hphantom{=\#\bigl\{}\text{both of whose sides are affinely spanned by }X\bigr\}\ \le\ 2^{k-1}-1 .\end{aligned}
\]
\end{proposition}

\begin{proof}
Write $f_1$ for the linear functional with $f_1|_H=\varphi$; then $f_2=f_1+f_0$ satisfies $f_2|_H=1+\varphi$. The supports of $w_{f_1},w_{f_2}$ are $X\cap\{\varphi=1\}$ and $X\cap\{\varphi=0\}$ — disjoint and covering $X$. By Lemma~\ref{lem:hyp} they are circuits iff $X\cap f_1^\perp=X\cap\{\varphi=0\}$ spans $f_1^\perp$ and $X\cap f_2^\perp=X\cap\{\varphi=1\}$ spans $f_2^\perp$. As $0\notin H$, a subset of $H\cap f^\perp$ spans the linear hyperplane $f^\perp$ iff it affinely spans the affine hyperplane $H\cap f^\perp$. Directions of parallel pairs correspond to nonzero linear parts, of which there are $2^{k-1}-1$; the pair $\{\varphi,1+\varphi\}$ describes one direction. The partition is determined by the direction.
\end{proof}

\src{New here (a reformulation; no statement from the release files is replaced). The equality $n_2=\#$directions was verified on all $231$ classified classes (Appendix~\ref{app:repro}).}

Thus W2 asks that \emph{two different directions of $\AG(k-1,2)$ cut $X$ into two affinely spanning halves}. The extreme case is $X=\AG(k-1,2)$, for which every direction works ($n_2=2^{k-1}-1=p-1$, Theorem~\ref{thm:L6}).

\subsection{Graver bases, lattice ideals and zonotopal tilings}

Graver bases of integer lattices are classical in integer programming and algebraic statistics \cite{Graver1975,Schrijver1986,Sturmfels1996}. The relation lattice $\Lambda_M\supseteq2\Z^n$ is a very special lattice for which the Graver basis is explicitly described by Proposition~\ref{prop:graver}: sign vectors on binary circuits plus the vectors $\pm2e_i$. A cell is a parallelepiped with Graver generators, i.e.\ a zonotope with independent generators \cite{Ziegler1995}; unions of such cells, and the question of when they fit together face to face, resemble the theory of zonotopal tilings, in which tiles of a tiling of a zonotope meet face to face by definition. In the present setting, however, we consider the set of \emph{all} admissible cells rather than a chosen tiling, and the question is whether this set is a complex; a collision is exactly a place where two cells of the family cannot belong to a common tiling in which they appear as tiles. Margolis, Saliola and Steinberg \cite{MSS} state the Bohne--Dress theorem as a bijection between zonotopal tilings of the zonotope $Z(\mathcal A)$ of a vector configuration and single-element liftings of the associated oriented matroid, with a tiling defined by four axioms: (ZT1) the tiles cover $Z(\mathcal A)$; (ZT2) faces of tiles are tiles; (ZT3) two tiles meet in a common face or not at all; (ZT4) every edge of $Z(\mathcal A)$ is a tile. Our face-to-face condition (Proposition~\ref{prop:facetoface}) is the analogue of (ZT3). The remaining axioms have no counterpart in our setting: our cells lie in the cube $[0,1]^n$ and are not required to cover a zonotope or to include its edges, so the theorem is not applicable as stated and we draw no conclusion from the analogy. We do not pursue the connection.

\subsection{Vector-space partitions}

A \emph{vector-space partition} of a finite vector space $V$ is a collection of subspaces such that every nonzero vector lies in exactly one of them; it has \emph{type} $[d_1^{n_1}\dots d_t^{n_t}]$ if it has $n_i$ members of dimension $d_i$ \cite{Heden2012}. The set $\{V_\Part/\langle\one\rangle\}$ of Theorem~\ref{thm:vsp} is a vector-space partition of $C/\langle\one\rangle\cong\F_2^{k-1}$, with parts of dimension $m_\Part-1$; its $1$-dimensional parts are the $n_2$ two-block partitions. What follows separates what is classical about this object from what is specific to the present setting.

\paragraph{Classical: the packing condition.} Heden's survey calls the identity $\sum_i n_i(q^{d_i}-1)=q^n-1$ the \emph{packing condition} (Section~5.1; it is not numbered in the survey). In our quotient it reads $\sum_\Part(2^{m_\Part-1}-1)=2^{k-1}-1$, which is the counting identity of Theorem~\ref{thm:counting} divided by $2$. So Theorem~\ref{thm:counting} is the classical packing condition; the content of Theorem~\ref{thm:vsp} is that the partition arises from circuit partitions of $\one$, not the identity itself. The \emph{dimension condition} $d_i+d_j\le n$ (attributed in the survey to Bu) reads $(m_\Part-1)+(m_{\Part'}-1)\le k-1$ here and is already a consequence of $V_\Part\cap V_{\Part'}=\{0,\one\}$.

\paragraph{Classical: types of small vector-space partitions.} According to the survey (Section~5.3), all types of vector-space partitions of $V(n,2)$ were determined for $n\le7$ by El-Zanati and coauthors, and Kurz \cite{KurzGF28} completes $\mathrm{GF}(2)^8$. The partitions arising here for $k\le6$ live in $\F_2^{k-1}$ with $k-1\le5$, so their types fall in the range where the list of admissible types is known. We did \emph{not} enumerate that list, and therefore did not compare the types of Appendix~\ref{app:J2} with it beyond the necessary condition checked below; the comparison remains open.

\paragraph{Classical: Heden's tail theorem, and a consequence derived here.} The \emph{tail} of a vector-space partition is the set of its parts of smallest dimension. Heden's theorem on the length $u_1$ of the tail, in the form quoted by Kurz \cite{KurzTail} (the source text of that statement is partly garbled; the cases below are as reconstructed there and in the survey): if the type is $d_l^{u_l}\dots d_2^{u_2}d_1^{u_1}$ with $u_1,u_2>0$, then (i) if $q^{d_2-d_1}\nmid u_1$ and $d_2<2d_1$ then $u_1\ge q^{d_1}+1$; (ii) if $q^{d_2-d_1}\nmid u_1$ and $d_2\ge2d_1$ then either $d_1\mid d_2$ and $u_1=[d_2]_q/[d_1]_q$, or $u_1>2q^{d_2-d_1}$; (iii) if $q^{d_2-d_1}\mid u_1$ and $d_2<2d_1$ then $u_1\ge q^{d_2}-q^{d_1}+q^{d_2-d_1}$; (iv) if $q^{d_2-d_1}\mid u_1$ and $d_2\ge2d_1$ then $u_1\ge q^{d_2}$.

\begin{proposition}[A consequence of Heden's theorem; derived here]\label{prop:heden-n2}
Let $C$ be a minimal non-embedded box with $n_2\ge1$ and with at least one partition of $\one$ having $m_\Part\ge3$ blocks, and let $d=\min\{m_\Part:m_\Part\ge3\}-1\ (\ge2)$. If Heden's theorem holds as quoted above, then $n_2\ge2^{d}-1\ge3$. In particular $n_2\ge2$ (the statement W2) holds in every such box.
\end{proposition}

\begin{proof}
In the quotient $F_2^{k-1}$ the partition has tail of dimension $d_1=1$ and length $u_1=n_2$, and the next smallest dimension is $d_2=d\ge2$ with $u_2>0$. As $d_2\ge2=2d_1$ only (ii) and (iv) can apply, with $q=2$. In case (iv), $u_1\ge2^{d_2}$. In case (ii), either $u_1=[d_2]_2/[1]_2=2^{d_2}-1$ or $u_1>2\cdot2^{d_2-1}=2^{d_2}$. In all cases $u_1\ge2^{d_2}-1\ge3$.
\end{proof}

\src{Derived here from the theorem as quoted in \cite{KurzTail}; the theorem itself is not re-proved. The inequality $n_2\ge2^{d}-1$ was checked on all $231$ classified classes ($87$ of them have a part with $m_\Part\ge3$), with no violation (\texttt{scripts/lit\_checks.py}). Status: \textsc{derived from a cited theorem}; not used anywhere in the proofs of this paper.}

\paragraph{What the classical theory does not give.} Proposition~\ref{prop:heden-n2} \emph{assumes} $n_2\ge1$. Nothing in the packing, dimension or tail conditions forces a $1$-dimensional part to exist (a tail may have larger dimension), and nothing relates $n_2$ to the number $p$ of reduced coordinates. The existence of a two-block partition (W2) and the bound $n_2\ge p-1$ for $p\ge9$ with its equality cases (Theorem~\ref{thm:L6}) use the weight function: the second-moment identity (Lemma~\ref{lem:moment}) and the Parseval identity of Theorem~\ref{thm:parseval}, which have no counterpart in the survey as read. In the other direction, the bound $n_2\ge p-1$ is much stronger than the classical $n_2\ge2^d-1$ when $p$ is large relative to $2^d$. With the caveat that we read only one survey and two papers on the subject, we found no statement in them that implies W2 for the partitions arising here.

\paragraph{Examples.} For $E_4$ ($k=3$) the quotient is $\F_2^2$ and the partition is $[1^3]$ (the three points). For $k=4,p=6$ (Appendix~\ref{app:G}) the quotient is $\F_2^3$ and the type is $[2^1\,1^4]$: a line and four points, in accordance with $u_1=4\ge2^2-1$. For $X=\AG(k-1,2)$ the type is $[1^{2^{k-1}-1}]$, the partition of $\F_2^{k-1}$ into points.

\subsection{Fourier analysis on \texorpdfstring{$\F_2^k$}{F2\^{}k}}

The weight enumerator of a binary linear code is determined by the Fourier transform of the indicator of its point set $X$ \cite{MacWilliamsSloane1977,Terras1999,ODonnell2014}. The Pless power-moment identities (stated for example as Theorem~5.1 in \cite{Kim2009}, where they are attributed to MacWilliams and Sloane) express the moments $\sum_jj^hB_j$ of the weight distribution of a code through the weight distribution of its dual; for $h=2$ and a dual of minimum distance at least $3$ they give the closed form of Lemma~\ref{lem:moment}. We proved Lemma~\ref{lem:moment} directly and did not verify the general identity beyond this reading. Lemma~\ref{lem:moment} is the second power moment; Theorem~\ref{thm:parseval} decomposes it along the vector-space partition. A natural question is whether higher moments yield further constraints (Question~\ref{q:moments}).

\subsection{Embedding obstructions}

The word ``embedding'' is used literally: let $K_{\mathrm{abs}}$ be the abstract cell complex whose cells are the $Z(S,x)$, with the face relations of Proposition~\ref{prop:facetoface}, and let $\iota:K_{\mathrm{abs}}\to\R^n$ be the realization. A point of $K_{\mathrm{abs}}$ lies in the relative interior of exactly one cell, so $\iota$ is injective iff the images of the relative interiors of distinct cells are disjoint iff there is no collision; thus the box is embedded iff $\iota$ is an embedding. Classical obstructions to embedding cell complexes in Euclidean space---for instance the van Kampen--Flores theorem \cite{Matousek2003}---are \emph{topological} (they apply to the abstract complex regardless of realization). The obstruction here is of a different nature: $K_{\mathrm{abs}}$ always exists, and non-embeddedness is a property of one particular map, forced by a purely linear-algebraic feature of the code (Theorem~\ref{thm:E1}). Whether $K_{\mathrm{abs}}$ could be embedded in some other way in $\R^n$ or in a higher-dimensional space is a separate question that we do not consider.

\subsection{What the comparison suggests}

The pieces that connect to standard areas are: minimal codewords and their disjoint decompositions (codes, matroids); hyperplane complements and affine subsets (finite geometry); vector-space partitions; Fourier analysis on $\F_2^k$. The pieces that are specific to the setting are the Graver cells, collisions and the witness rank, which have no direct standard counterpart that we know.

%% -------------------------------------------------------------------
\section{Provenance, verification and reproducibility}\label{sec:prov}

\subsection{Tiers of evidence}

\begin{enumerate}[label=\textbf{T\arabic*.},leftmargin=*]
\item \emph{Source-reported proof.} A proof recorded in the GEO-I releases (v006--v011; Section~\ref{sec:intro}). Statements and, where available, proofs were read; none was independently verified by an outside party.
\item \emph{Internal written audit.} The re-derivation of the arguments for this manuscript by its authors (the model-assisted drafting process described in the acknowledgements). A second reading by the same hand; not an external review.
\item \emph{Independent computational tests.} Programs written afresh for this manuscript (Appendix~\ref{app:repro}) that test finite instances: exact integer / bitmask arithmetic throughout (no floating point is used in the mathematical checks).
\item \emph{External review.} None has taken place.
\item \emph{Formal proof.} None.
\end{enumerate}

Each universal theorem of Parts~II--III is proved in the text (tiers T1+T2) with one exception: the case $p\le8$ of Theorem~\ref{thm:W2} rests on an exhaustive computer enumeration (Proposition~\ref{prop:small}, tier~T3 in the sense of being a complete finite verification, with a different route available in the source).

\subsection{Finite computations accompanying the paper}

\begin{table}[ht]
\centering\small
\begin{tabular}{@{}p{5.3cm}p{4.3cm}p{3.0cm}p{1.7cm}@{}}
\toprule
Computation & Range & Instances & Failures\\
\midrule
E1 criterion vs.\ direct cell/collision enumeration; $\rmin=3$ & all multisets over $\F_2^{2},\F_2^3\setminus0$, $n\le7$, plus random codes over $\F_2^4\setminus0$ with $n\in[5,8]$ (up to $1500$) & $4213$ codes ($706$ non-embedded) & $0$\\
Circuit test: minimal supports vs.\ hyperplane rank test & same codes & $4213$ & $0$\\
Classification by exhaustive enumeration & $k=3,4,5$ & $1+93+57\,541$ spanning subsets & ---\\
Orbit counts / Burnside & $\AGL(m,2)$, $m=2,3,4$ & $24,1344,322\,560$ elements & ---\\
Identities (counting, Parseval, vector-space partition, pairwise $V_P\cap V_Q=\{0,\one\}$) & all minimal classes $k\le6$ & $16+215$ classes & $0$\\
Centre cell count (Proposition~\ref{prop:centre-count}) & all $16$ classes with $k\le5$ & $16$ & $0$\\
Two-block partitions = hyperplane directions (Proposition~\ref{prop:geo-n2}) & all $231$ classes & $231$ & $0$\\
Finite check for $p\le8$ (Proposition~\ref{prop:small}) & $3\le k<p\le8$ & $655$ candidates & ---\\
$215$ nullity-6 classes pairwise inequivalent & invariants; exhaustive search for $8$ pairs & $215$ & $0$ identifications\\
Heden tail bound $n_2\ge2^d-1$ (Proposition~\ref{prop:heden-n2}) & all $231$ classes & $87$ with a part $m_\Part\ge3$ & $0$\\
Cycle codes of multigraphs: E1 vs.\ Cooper--Okur cactus criterion (Remark~\ref{rem:cooperokur}) & random multigraphs, $\le5$ vertices, $\le9$ edges & $2000$ & $0$ mismatches\\
Source checkers (unchanged) on source certificate files & \texttt{checker11} (PARTITION), \texttt{checker10} (RANK6), \texttt{checker8} (RANK) & $231,215,12$ & $0$\\
Random spanning subsets non-embedded (Question~\ref{q:spanning}) & $k\in\{6,7,8\}$ & $540$ & $0$ embedded\\
\bottomrule
\end{tabular}
\caption{Finite computations performed for this manuscript. They are support for statements and examples, never proofs of universal claims (with the exception of the finite enumeration of Proposition~\ref{prop:small}, which is a complete verification of a finite statement).}\label{tab:checks}
\end{table}

\subsection{Source-reported items not re-derived}

The following were taken from the releases without independent derivation here: the completeness of the nullity-$6$ list (\gid{G10-T4}: $381$ orbits, $350$ spanning, $349$ non-embedded, $215$ minimal); the Burnside count $382$ for $m=5$; the supplementary scans \gid{G11-V1} (knapsack feasibility), the random nullity-$7$/$8$ tests (\gid{v11\_w2.json}: $154$ and $2$ boxes, no violation); and the statement \gid{G10-K} for $K_{2,t}$ ($t\le10$).

\subsection{Disclosures}

\begin{enumerate}
\item The GEO-I releases v001--v012 are cited at the fixed commit \geoicommit{} of the repository. The drafting process could not itself fetch the repository and read the local release folders supplied by the author instead; the local copies were not compared byte for byte with the pinned commit. No frozen release and no repository file was modified.
\item Of the bibliography entries other than the GEO-I releases, six were read in full text for Section~\ref{sec:literature} (\cite{Heden2012,KurzTail,KurzGF28,CooperOkur,MSS,Kim2009}); the others were assembled from catalogue records or author and publisher pages and were not read. Section~\ref{sec:literature} states which is which.
\item The figures are generated by scripts in the package from exact data; captions were written by hand and checked against the data. Three-dimensional renderings are exact pictures of the $3$-dimensional objects they show (the $E_3$ triangle, the octahedron in the hyperplane $\sum x_i=2$) and are never presented as renderings of higher-dimensional geometry.
\item Theorem~\ref{thm:rank3} depends on the enumeration of Proposition~\ref{prop:small}; all other steps are hand proofs.
\item Items marked ``new here'' (Appendix~\ref{app:crosswalk}) were derived in the course of writing this manuscript and were not found in the release files read; they have not been refereed.
\end{enumerate}

\paragraph{Acknowledgements.} The drafting, proof re-derivations, scripts and figures of this manuscript were produced in a collaboration between the author and an AI assistant (Claude, Anthropic). The mathematical programme, the source releases and all decisions on content are the author's. No external funding or external review is acknowledged.

%% -------------------------------------------------------------------
\section{Open questions}\label{sec:open}

\begin{question}[Beyond the $0/1$ box]\label{q:beyond}
What happens for boxes $Q>\one$ in binary relation lattices (cells with corners at coordinates up to $Q_i$), where Graver vectors $\pm2e_i$ now matter? Is there an analogue of Theorem~\ref{thm:E1}, and is $\rmin$ still uniformly bounded?
\end{question}

\begin{question}[Other lattices and moduli]\label{q:moduli}
For congruence lattices modulo $q>2$, or with several residue conditions, circuits are replaced by the (much larger) Graver basis. Is there an exact criterion for collisions of cells in terms of decompositions of lattice vectors into Graver vectors? For which lattices is the minimum witness rank uniformly bounded?
\end{question}

\begin{question}[Topology versus embedding]\label{q:topology}
The realized union $\iota(K_{\mathrm{abs}})$ of a non-embedded box has homology that need not be that of $K_{\mathrm{abs}}$. How do the Betti numbers of the image relate to the collisions? For $E_4$ the three squares sharing diagonals give a specific small complex; is there a general formula in terms of the partition data?
\end{question}

\begin{question}[Collision structures]\label{q:collisions}
Classify, for a given box, the set of occurring witness ranks and cell-dimension pairs. For $E_6$ the pairs are $(2,2)$ and $(3,3)$ with ranks $3$--$6$; what is the general pattern? Is the number of rank-three collisions determined by $(k,p,\text{type})$?
\end{question}

\begin{question}[Types of vector-space partitions]\label{q:types}
Which types of vector-space partitions of $\F_2^{k-1}$ arise as $\{\bar V_P\}$? In the classified range there are eight types for $k=6$ (maximum block number $5$). Is the number of blocks of a partition of $\one$ bounded in terms of $k$ more strongly than by $m\le p/2$? Can the minimal boxes with $n_2=p-1$ be classified for $p<9$?
\end{question}

\begin{question}[Higher moments]\label{q:moments}
Do the third and fourth moments of the weight distribution, decomposed along the partition, give further constraints that prove structural statements (for instance an upper bound on $m_P$) without enumeration?
\end{question}

\begin{question}[Algorithmic complexity]\label{q:algo}
What is the complexity of deciding embeddedness of a given code? Corollary~\ref{cor:cert} puts the problem in NP in the binary-matrix encoding; is it in P, or NP-complete? Is it fixed-parameter tractable in the nullity $k$?
\end{question}

\begin{question}[Spanning subsets]\label{q:spanning2}
Is every spanning subset $X\subseteq\AG(k-1,2)$ with $|X|\ge k+1$ non-embedded (Question~\ref{q:spanning})? A positive answer would give a complete characterisation of non-embedded boxes with $\one\in C$ in finite-geometric terms.
\end{question}

\begin{question}[Classification beyond $k=6$]\label{q:class7}
Re-derive the nullity-$6$ classification independently (for instance by the simplex normal form of Lemma~\ref{lem:simplex} and orbit enumeration), and determine the number of minimal classes for $k=7$.
\end{question}

\begin{question}[Cut spaces]\label{q:cutspaces}
For which graphs is the cut space a minimal box? Is $K_{2,t}$ minimal for all $t\ge3$ (checked for $t\le10$ in the source), and is its $\rmin$ explained by a family of two-bond partitions?
\end{question}

\begin{question}[Formal verification]\label{q:formal}
The proofs of Parts~II--III are short and elementary; they are natural targets for formalization in a proof assistant, with the finite enumeration for $p\le8$ replaced by a verified computation.
\end{question}
